% Korean edition of sections/s4_4_validity.tex.
\seclabel
\claim{\secnum{4.4}{sec:validity}\quad 계정 하나로 오전 한나절 한 테스트라는 점이 결과를 제한하며, 질문 5개가 남아 있다}

\lead{3\,s 폴링을 알려준 2026년 10월 28일 브라우저 기록을 빼면, 이 보고서의 모든 측정은 계정 하나로 2026년 11월 24일 05:38\textasciitilde 10:34 KST 오전 한나절에 이루어졌다.} 모든 요청은 포트폴리오 하나를 썼고, 한 네트워크의 PC 한 대에서 웹 클라이언트와 같은 경로·헤더·본문으로 curl과 래퍼의 httpx 클라이언트를 통해 보냈다. 따라서 부하, 날짜, 네트워크에 따라 달라지는 값은 다른 날 다를 수 있다. \ref{sec:orders}절의 체결 속도, \ref{sec:feed}절의 응답 크기, 그리고 모든 시간 측정이 그렇다. 웹 클라이언트가 보내지 않는 형태의 요청이나 두 번째 계정은 시험하지 않았다.

\lead{근거는 실패 비용이 가장 큰 곳에서 가장 얇다.} 래퍼가 넣은 주문을 피드에서 찾는 15\,s 시간은 \code{AddOrder} 후 1\textasciitilde 3\,s 만에 찾은 래퍼의 실제 매수와, 08:33까지 피드에 나타난 08:30 주문에 근거한다. 매도의 시간은 재지 않았고 느린 서버도 시험하지 않았다. 이 시간이 너무 짧으면 넣은 주문이 \code{AmbiguousState}가 된다. 반면 틀려도 지연만 생기는 연결 시간 측정은 요청 50개(5개씩 10회, \ref{sec:transport}절)에 근거한다.

\lead{계층별 시간은 대부분 09:00 개장 전에 쟀으므로, 장중 부하는 Python 클라이언트에 대해서만 측정됐다.} \ref{sec:transport}절의 16.0\ms{}와 106.5\ms{}는 개장 전 curl에서, 재사용 연결에서 래퍼의 31\textasciitilde 39\ms{}는 장중에 나왔다. 두 값은 클라이언트, 시각, 부하가 한꺼번에 다르며 보고서는 이 원인들을 분리하지 않는다.

\lead{유휴 타임아웃, 영업일 전환, 토큰 허용치, \code{Logout}, 피드에 주문이 나타나는 시간, 401 뒤 재전송은 관찰 2회 이하에 근거한다.} 15\textasciitilde 30\,s 유휴 타임아웃은 5\,s와 15\,s 침묵 뒤 재사용된 연결 1개와 30\,s 뒤 교체된 연결 1개에 근거한다. 06:02:10\textasciitilde 07:06:44 사이의 영업일 전환은 하루 동안의 \code{currBizDay} 두 번 읽기에 근거한다. 토큰 만료 후 300\,s 허용치는 테스트 2회(\ref{sec:auth}절)에, 저장한 쿠키 묶음이 \code{Logout} 뒤에도 갱신된다는 결과는 저장한 쿠키 묶음 1개로 한 호출 1회에 근거한다. \code{AddOrder}에서 피드에 주문이 나타나기까지는 래퍼의 실제 매수에서 1\textasciitilde 3\,s였고, 그 밖에는 상한으로만 알려져 있다(08:30:05 주문은 08:33까지 피드에 있었다). 따라서 \ref{sec:orders}절의 15\,s 검색 시간은 느린 서버에서 확인하지 않았다. 401 뒤 주문 재전송은 주문 요청에서 관찰한 적이 없고 프레임워크 문서에만 근거한다(\ref{sec:auth}절). 래퍼의 매개변수가 이 중 하나에 의존하는 곳에서는 관찰값 안쪽으로 여유를 두었다(부록~\ref{app:decisions}). 어느 것도 다른 날 다시 측정하지 않았다.

\lead{서버의 동작은 문서화되어 있지 않으므로 이 결과들은 예고 없이 바뀔 수 있다.} 델타 규칙, \code{wei}의 의미, \code{qty2ord}의 내림, 300\,s 허용치는 응답에서 추론했으며 주최측이 확인해 주지 않았다. 래퍼는 확인할 수 있는 것은 실행 중에 검사한다. \code{currBizDay}가 KST 날짜와 다르면 주문을 거부하고, 마지막 피드 응답이 26\,s보다 오래되면 \code{StaleFeed}를, 넣은 주문을 찾지 못하면 \code{AmbiguousState}를 던진다. \code{tests/live}의 실서버 검사는 측정을 반복한다. 그중 3개(\code{live_orders}, \code{live_sell}, \code{live_rounding})는 실제 주문을 넣고 정지한다.

\lead{주최측 한도는 IP 주소 단위이므로, 공유 NAT 주소 하나 뒤에 있는 동아리 회원들은 한도를 함께 쓴다.} 주최측 답장[5]은 하루 100건 이상 주문하는 IP 주소와, 시스템에 큰 부하를 주는 트래픽 폭주가 차단될 수 있다고 밝힌다. 래퍼의 하루 80건 한도는 계정 1개만 센다. 같은 캠퍼스 네트워크에서 래퍼를 돌리는 회원들은 같은 공인 주소로 나가므로, 이들의 주문과 요청이 래퍼가 볼 수 없는 한도에 함께 계산될 수 있다. 부하 기준 자체는 알려지지 않았다. 시간 측정은 순차 요청 5개씩 보냈고 부하 테스트는 하지 않았다.

\lead{보장 2개는 운영 가정 아래에서만, 또는 한 시점에서만 성립한다.} 첫째, 주문 응답에 ID가 없으므로 래퍼는 같은 순간 웹 페이지나 다른 프로세스에서 넣은 똑같은 주문과 자기 주문을 구분할 수 없다. 맞는 새 레코드가 둘 이상이면 \code{AmbiguousState}를 던지지만 외부 주문 하나와 맞는 경우는 알아챌 수 없으므로, 래퍼는 실행 중에 계정당 주문 제출자가 하나라고 가정한다. 둘째, \code{plan()}은 계획 시점의 데이터로 목표 비중을 종목·섹터·합계 한도와 대조한다. 매도가 매수보다 먼저 나가지만 먼저 체결된다는 보장은 없으므로(7주가 7분 37초에 걸쳐 체결됨, \ref{sec:orders}절), 리밸런싱이 체결되는 동안 섹터나 합계가 한도를 넘을 수 있다. 주최측이 위반을 언제 재는지는 질문 Q4다. \figref{guarantees}는 사용자가 믿을 수 있는 것, 각 보장에 필요한 가정, 근거, 시험하지 않은 것을 나열한다.

\lead{실서버에서 일부러 장애를 일으키지는 않았고, 실서버 테스트 뒤에 바뀐 코드는 오프라인에서만 돌았다.} 실서버를 상대로 피드 끊김, 재연결, 백오프, \code{AddOrder} 도중 연결 끊김을 일으키지 않았고, 한 실행 안에서 720\,s를 넘겨 갱신한 프로세스도 아직 없다. 실서버 테스트 뒤 동작 6가지를 바꿨다. 피드 요청마다 응답 마감 시간 하나, 저장한 쿠키 만료 시각, 재전송 실패 뒤 검색, 이미 돌려준 주문 ID 제외, 주문 하나에 새 레코드가 둘 이상 맞을 때 \code{AmbiguousState}, 어떤 요청에서든 403이나 429 뒤 쓰기 중단이다. 오프라인 테스트 31개 중 7개가 이 변경을 다루며, 어느 것도 실서버에서 돌지 않았다.

\lead{비밀번호를 바꾼 뒤에도 저장된 쿠키 묶음이 갱신되는지는 시험하지 않았다(\figref{openq}의 Q3).} 시험할 때까지 래퍼는 \code{session.json}을 \code{Logout} 뒤에도 갱신되는 7일짜리 자격 증명으로 다루며(\ref{sec:auth}절), 저장소 밖에 권한 0600으로 보관한다(\figref{security}).

\lead{주문을 자동화하고 웹 클라이언트의 요청을 재현해도 된다는 근거는 2026년 11월 24일 주최측 이메일[5] 하나뿐이다.} 답장은 주최측이 공식 API를 제공하지 않으며, 스크립트나 웹 드라이버로 주문을 자동화하는 것을 금지하지 않는다고 밝힌다(\ref{sec:intro}장). 사본은 보관하고 있다. 이 입장은 우리 문의에 대한 답장에만 있으며, 주최측은 언제든 바꿀 수 있다.

\lead{질문 5개가 남아 있다. 4개는 테스트로, 1개는 주최측만 답할 수 있다(\figref{openq}).} 질문마다 래퍼는 답이 반대로 나와도 안전한 잠정 동작을 따른다. 예를 들어 취소가 100건에 포함되는지 주최측이 답할 때까지 정지를 주문으로 센다(Q2). 이전 질문 3개는 11월 24일에 해결됐다. \code{qty2ord}는 내림된다. 1.76주어치 삼성전자 주문이 두 번 모두 1주를 요구했다(\ref{sec:oneshare}절). 주문 직후 \code{Accepted} 상태를 관찰했으므로(\ref{sec:orders}절) Q5는 이제 \code{Generated}와 \code{Deleted}, 그리고 체결되지 않은 채 \code{hm1}에 이른 시간대 주문이 어떻게 끝나는지를 묻는다. \code{bs_updated}는 모든 델타에서 잔고 필드 \code{prft}, \code{na}, \code{rt}, \code{nav}, \code{stkAmt}를 싣고 온다(\ref{sec:feed}절).

% Korean edition of figs/openq.tex.
\fig[wide]{openq}{질문 5개가 남아 있다. 4개는 테스트로, 1개는 주최측만 답할 수 있다}[미해결 질문, 해결될 때까지 래퍼가 하는 일, 다음 단계]{%
{\tabfont\begin{tabularx}{\linewidth}{>{\RaggedRight\arraybackslash}p{6mm} >{\RaggedRight\arraybackslash}X >{\RaggedRight\arraybackslash}p{50mm} >{\RaggedRight\arraybackslash}p{50mm}}
\thd{} & \thd{질문} & \thd{잠정 동작} & \thd{해결 방법(다음 단계)}\\
Q1 & 클라이언트 코드에서 가져온 \code{CancelOrder}, \code{ModifyOrder}, \code{DeleteOrder} 본문이 동작하는가? 어느 것도 실서버에서 시험하지 않았다 & 래퍼는 쓰지 않는다. 주문은 \code{CancelOrderBatch}로 정지한다 & 웹 페이지가 어떤 주문에 이 기능을 제공할 때 요청을 캡처한다\\\rowsep
Q2 & 취소가 100건 한도에 포함되는가? & 정지한 주문마다 1건으로 센다. 래퍼는 하루 80건으로 제한한다 & 주최측에 문의한다\\\rowsep
Q3 & 비밀번호를 바꾼 뒤에도 저장한 쿠키 묶음이 갱신되는가? & 쿠키 묶음을 7일짜리 자격 증명으로 다룬다(\ref{sec:design}) & 비밀번호를 바꾼 뒤, 바꾸기 전에 저장한 쿠키 묶음으로 갱신해 본다\\\rowsep
Q4 & 서버는 섹터·소형주 위반을 언제 세는가? 장중 아무 때나인가, 장 마감인가? 섹터 한도는 시장 비중을 따라 움직이고, 리밸런싱 중에는 매도가 매수보다 먼저 체결된다는 보장이 없어 비중이 잠시 한도를 넘을 수 있다(\ref{sec:guidelines}) & \code{plan()}은 호출 때마다 피드에서 읽은 한도로 계획 시점의 목표 비중을 검사한다. \code{guidelines()}는 소형주 합계를 보고한다 & 매시간 \code{Violations}와 피드의 섹터 비중을 기록해 비교한다\\\rowsep
Q5 & 주문 상태 \code{Generated}와 \code{Deleted}는 언제 나타나며, 체결되지 않은 채 \code{hm1}에 이른 시간대 주문은 어떻게 끝나는가? & \code{Done}이 아닌 상태는 모두 진행 중으로 본다. \code{plan()}은 남은 크기를 빼고 \code{wait()}은 계속 기다린다. \code{cnclT}가 null인 \code{Done} 주문은 체결로 보고한다 & 서버가 거부한 주문을 포함해 \code{ords}의 모든 상태 변화를 기록한다\\
\end{tabularx}}}{11월 24일에 해결되어 이 목록에서 뺀 것: 실시간 피드가 밀어 보내는지(\ref{sec:feed}), \code{updated}의 형식과 델타 병합의 정확성(\ref{sec:feed}), 동시 갱신(\ref{sec:auth}), \code{Orders} 매개변수(\ref{sec:queries}), 체결·정지된 주문의 최종 상태와 주문 직후 \code{Accepted} 상태(\ref{sec:orders}), 매도 주문 형식(\ref{sec:wrappertests}), \code{qty2ord}의 반올림 방식(내림, \ref{sec:oneshare}), \code{bs_updated}의 내용(\ref{sec:feed}), 순자산의 출처(\ref{sec:oneshare}), 가이드라인의 근거가 되는 피드 필드(\ref{sec:guidelines}).}{DART Quant Team}

% Korean edition of figs/guarantees.tex.
\fig[wide]{guarantees}{사용자가 믿을 수 있는 것과 그 가정}[\code{pytimefolio} 0.1.0의 보장마다 필요한 가정, 근거, 시험하지 않은 것]{%
{\tabfont\begin{tabularx}{\linewidth}{>{\RaggedRight\arraybackslash}p{38mm} >{\RaggedRight\arraybackslash}p{42mm} >{\RaggedRight\arraybackslash}X >{\RaggedRight\arraybackslash}p{44mm}}
\thd{사용자가 믿을 수 있는 보장} & \thd{필요한 가정} & \thd{근거} & \thd{시험하지 않은 것}\\
병합 상태는 \code{t}가 같은 서버 스냅샷과 같다 & 피드 요청은 하나만. 서버가 델타 규칙 (a)를 유지 & 실서버: 09:35에 \code{pos}와 \code{ords}가 15회 중 15회 일치(\ref{sec:feed}) & \code{bs}와 \code{msecs}의 필드 단위 비교. 장중 규칙 (b)와 (c)\\\rowsep
읽기, 주문, \code{wait()}은 26\,s보다 오래된 상태를 쓰지 않는다 & 피드 태스크가 실행 중. 3\,s마다 폴링은 목표이지 보장이 아님 & 오프라인: 낡은 피드는 \code{StaleFeed}. 핑은 20\,s 응답 마감을 늘리지 않음(\ref{sec:wrappertests}) & 실제 피드 장애\\\rowsep
\hl 서버가 받았을 수 있는 주문은 다시 보내지 않는다 & 401은 주문 처리기보다 먼저 온다[9]. 연결 설정 전 연결 오류는 요청이 나가지 않았음을 뜻함 & 오프라인: 잃어버린 응답은 재전송 안 함. 실패한 재전송은 검색하며 세 번째 전송 없음(\ref{sec:orders}) & 주문에서 받은 401. 실제 \code{AddOrder} 도중 연결 끊김\\\rowsep
\hl 넣은 주문은 15\,s 안에 확인되거나 \code{AmbiguousState}가 발생한다 & 서버가 15\,s 안에 새 주문을 피드나 \code{Orders}에 보여줌 & 실서버: 매수는 \code{AddOrder} 후 1\textasciitilde 3\,s 만에 찾음. 08:30 주문은 08:33까지(\ref{sec:validity}) & 느린 서버. 매도의 지연\\\rowsep
\hl 돌려준 주문 ID는 사용자의 주문이다 & 래퍼 실행 중 계정당 제출자 하나 & 오프라인: 같은 주문 2건은 ID 2개. 맞는 새 레코드 2개면 \code{AmbiguousState}(\ref{sec:orders}) & 웹 페이지나 다른 프로세스에서 넣은 똑같은 주문 1건은 알아챌 수 없음\\\rowsep
하루 최대 80건, 쓰기 사이 최소 1\,s & 건수 파일 하나에 프로세스 하나. 취소를 주문으로 셈 & 오프라인: 한도, 동기화, 1.002\,s 간격. 실서버: 10:34까지 카운터 80건 중 16건(\ref{sec:design}) & 건수 파일 하나를 쓰는 여러 프로세스. 주최측이 취소를 세는지(Q2)\\\rowsep
\hl 목표 비중은 계획 시점에 종목, 섹터, 합계 한도를 지킨다 & \code{plan()} 실행 시점에 피드 값(\code{weiLimit}, \code{msecs})이 최신 & 실서버: 산업재 +6\%p 거부, +3\%p 허용. 100\% 초과 목표 거부(\ref{sec:guidelines}) & 실행 중 비중: 먼저 보낸 매도가 먼저 체결된다는 보장 없음. 위반을 언제 세는지(Q4)\\\rowsep
403이나 429 뒤에는 모든 쓰기가 멈춘다 & 서버가 차단을 403이나 429로 알림 & 오프라인: 쓰기나 읽기에서 403·429면 쓰기 중단. 차단된 피드는 재연결 중단(\ref{sec:design}) & 실서버에서 차단을 본 적 없음\\\rowsep
최대 7일간 비밀번호 없이 갱신 & 서버가 7일 쿠키 만료를 유지. 720\,s에 갱신 & 실서버: 두 번째 프로세스가 저장된 쿠키 묶음으로 갱신. 오프라인: 재시작해도 만료 유지(\ref{sec:auth}) & 720\,s 이상 또는 7일 동안 돈 프로세스\\
\end{tabularx}}}{강조 행은 운영 가정에 기대거나, 실패 비용에 비해 근거가 얇은 보장이다. 괄호 안 절 번호가 근거다. \figref{goalstatus}는 같은 테스트로 목표를 판단한다.}{DART Quant Team; \figref{tests}}

